\chapter{Series of Functions and the Exchange of Limits}

A small numerical library computes $\ln(1+x)$ with the series
$\sum_{k\ge 1}(-1)^{k+1}x^{k}/k$. Its author tested the function at $x=0.5$
with twelve terms and compared it against a standard library: the error was
$6.4\times 10^{-6}$. The test passed, the documentation states that the series is
accurate to six digits, and the code went into production. The first user called it with
$x=0.99$. Twelve terms give $0.6529204447$, whereas
$\ln 1.99=0.6881346387$. The error is $3.5\times10^{-2}$, about five thousand times
larger, for the same function with the same number of terms.

That damage is not a typo. For every $x$ in $(-1,1)$ the series converges
to $\ln(1+x)$; what changes is only how many terms are needed: $16$
terms at $x=0.5$, $89$ at $x=0.9$, $721$ at $x=0.99$, $5248$ at $x=0.999$,
all of them for an error below $10^{-6}$. No single number of terms suffices
for all $x$, because $x$ may lie as close as we please to $1$.

Compare with a closed subinterval $[-r,r]$, $0<r<1$. There one estimate
holds simultaneously for all $x$: the remainder of an alternating series is bounded by
$r^{N+1}/(N+1)$, and that bound does not mention $x$. For $r=0.5$ and $N=12$
the bound is $9.4\times10^{-6}$, the actual error $6.4\times10^{-6}$. This is the
difference that this entire chapter is about.

Chapter 5 dealt with series of numbers, and each $x$ gives a different series of
numbers. Chapters 11 and 13 define continuity and the integral point by point,
so neither of them has a way to move a limit past an integral sign without
extra conditions. The tool still missing is a distance between two functions; once that
distance exists, the whole family of series becomes a single sequence in a metric space.

\begin{untukanak}
Imagine a class with thirty students, and the teacher writes a single number
on the board: one hundred. Each student's task is to guess that number, and after every
guess the teacher announces whether the guess is getting closer. The first student is very
fast: within five guesses he is off by less than one. The second student is
slow: he needs twenty guesses. The third is very slow: he needs a hundred.
But all three, given enough time, end up off by less than one.

Now the teacher changes the question. Not ``does every student eventually get
close'', but ``is there a moment when all students are already close at the same
time''. For that class the answer is no: the slowest student is always
lagging, and one can always add a student who is slower still.

The same question arises for the formulas a computer uses. A formula
may be correct for every value put into it, yet the number of terms
needed can blow up without bound as the value approaches the edge. If a
program takes ``correct for every value'' to mean ``one formula suffices
for all values'', that program will be far off at the most troublesome values, and the
error will not show up in the test its author chose.

In this chapter we shall see why those two questions are not a matter of words,
but a matter of whether two operations may be exchanged in order: limit with
integral, limit with derivative, and summation with both.
\end{untukanak}

\section{Pointwise and uniform convergence}
\tingkat{basic}

Take a family of functions $f_1,f_2,\ldots$ all defined on a single
set $E\subseteq\mathbb{R}$. For one point $x$ that we hold fixed,
the values $f_1(x),f_2(x),\ldots$ are nothing but a sequence of numbers, and Chapter 4
already said what it means for that sequence to converge. Hence the first definition
almost writes itself.

\begin{definisi}
Suppose $f_n$ and $f$ are defined on $E\subseteq\mathbb{R}$.
\begin{enumerate}
\item We say $f_n\to f$ \emph{pointwise}\index{pointwise convergence} on $E$
      if for every $x\in E$ and every
      $\varepsilon>0$ there is $N$ such that $n\ge N$ implies
      $|f_n(x)-f(x)|<\varepsilon$. Here $N$ may depend on $x$ and
      on $\varepsilon$.
\item We say $f_n\to f$ \emph{uniformly}\index{uniform convergence} on
      $E$ if for every $\varepsilon>0$ there is $N$ such that $n\ge N$
      implies $|f_n(x)-f(x)|<\varepsilon$ for \emph{all} $x\in E$
      at once. Here $N$ depends only on $\varepsilon$.
\item If $h$ is bounded on $E$, the supremum norm\index{supremum norm} of $h$
      is $\|h\|_\infty=\sup_{x\in E}|h(x)|$.
\end{enumerate}
\end{definisi}

\kamus{uniform convergence}{konvergensi seragam}{all points must already be close
at the same time, not one by one on their own schedules.}

One quantifier changes place, and the whole content of this chapter hangs on
that move. In the first definition the order is ``for every point there is a
threshold''. In the second the order is ``there is one threshold for all points''.
The first sentence is true for the series in the chapter opening, the second is false, and
the difference between them is the difference between $16$ terms and $5248$ terms.

\begin{bacakeras}
Two voices take turns.

First voice: ``Does this series get there?''
Second voice: ``It gets there. Slowly, but it gets there.''
First voice: ``Does it get there for this point?''
Second voice: ``For this point. For the point next door, another time, at its own pace.''

That is the first promise. Every point may be slow, as long as it has one tempo
of its own. A point far from the edge arrives within a dozen steps;
a point pressed against the edge needs thousands.

Now the first voice asks in other words: ``Is there one tempo for
all points at once?''

The second voice stays silent longer. ``On a closed interval that stops before the edge,
yes. On the whole open interval, no. And what is missing there is not a failure of
counting; it merely demands a tempo that moves together with the point.''

Repeat this line until it sits on the tongue: all points may be slow, as long as one
tempo. If the tempo must differ for each point, its name is not uniform.
\end{bacakeras}

\begin{contoh}
Suppose $f_n(x)=x^n$ on $E=[0,1]$. For $0\le x<1$ we have $x^n\to0$,
whereas $x^n=1$ for all $n$ if $x=1$. So $f_n$ converges pointwise
to
\[
  f(x)=\begin{cases}0, & 0\le x<1,\\ 1, & x=1.\end{cases}
\]
That convergence is not uniform, since $\sup_{x\in[0,1]}|x^n-f(x)|=1$ for
every $n$: at $x=1$ the difference is exactly $1$, and for $x$ close enough to $1$
from the left $x^n$ approaches $1$ while $f(x)=0$.
\end{contoh}

That example shows two things. First, the pointwise limit of continuous functions
need not be continuous: every $x^n$ is continuous, its limit jumps at $x=1$. Second,
uniformity depends on the set, not only on the formula of the function.

That dependence can be measured. On $[0,r]$ with $r<1$ the convergence is
uniform, since $\sup_{[0,r]}|x^n|=r^n\to0$. What is interesting is when the right
endpoint also moves with $n$. On $[0,1-1/n]$ the supremum is
$(1-1/n)^n\to1/e=0.3678794412$; so shrinking by $1/n$ does not yet
rescue uniformity. On $[0,1-1/\sqrt{n}]$ the supremum is
$(1-1/\sqrt n)^n$, which collapses like $e^{-\sqrt n}$: $2.2\times10^{-2}$ at
$n=10$, $2.7\times10^{-5}$ at $n=100$, $2.2\times10^{-44}$ at $n=10^4$. The rate
of shrinking, not the formula of the function, is what decides.

The diagonal test yields sharper additional information. Take the points
$x_n=1-1/n$ chasing the break point. Then $f_n(x_n)=(1-1/n)^n\to1/e$,
while $f(x_n)=0$ for all $n$. So not only does the supremum fail to
shrink: the values at well-chosen points also refuse to approach the
pointwise limit. This pattern is used repeatedly in this chapter. To
refute uniformity, find one sequence of points that keeps the difference
large.

Two simple families separate the two symptoms. In the first, $g_n$ is a
triangle of height $1$ with base $1/n$ and apex at $1-1/n$: the apex
moves right, the height never drops, and the area is $1/(2n)$. A point
$x<1$ is touched by the apex only while $1/(2(1-x))<n<3/(2(1-x))$, so at $x=0.9$
only nine values of $n$ give a positive value, at $x=0.99$ ninety-nine,
and once the last apex has passed the value is zero forever. The limit is zero
at every point, whereas $\|g_n\|_\infty=1$. In the second, $h_n$ has height $1/n$
with base $1$ in the same place: the apex shrinks instead of moving, and
$\|h_n\|_\infty=1/n\to0$ (Figure~\ref{fig:puncak}).

\begin{figure}[htbp]
\centering
\begin{tikzpicture}[x=5.6cm,y=3.1cm,font=\small]
  \draw[->,layergrey] (0,0)--(1.15,0) node[right,font=\scriptsize]{$x$};
  \draw[->,layergrey] (0,0)--(0,1.32) node[above,font=\scriptsize]{$y$};
  \node[below,font=\scriptsize,layergrey] at (1,0) {$1$};
  \foreach \n/\c in {16/layerindigo!15,8/layerindigo!30,4/layerindigo!50,2/layerindigo!75}{
    \fill[\c] ({1-1/\n-1/(2*\n)},0) -- ({1-1/\n},1) -- ({1-1/\n+1/(2*\n)},0) -- cycle;
  }
  \draw[dashed,layergrey] (0,1)--(1.06,1);
  \node[above,font=\scriptsize,layergrey] at (0.72,1) {height $1$};
  \begin{scope}[xshift=7.0cm]
    \draw[->,layergrey] (0,0)--(1.15,0) node[right,font=\scriptsize]{$x$};
    \draw[->,layergrey] (0,0)--(0,1.32) node[above,font=\scriptsize]{$y$};
    \node[below,font=\scriptsize,layergrey] at (1,0) {$1$};
    \foreach \n/\c in {8/layerteal!18,4/layerteal!32,2/layerteal!52,1/layerteal!72}{
      \fill[\c] (0,0) -- (0.5,{1/\n}) -- (1,0) -- cycle;
    }
    \node[right,font=\scriptsize,layergrey] at (0.52,0.5) {$n=2$};
    \node[right,font=\scriptsize,layergrey] at (0.52,0.25) {$n=4$};
    \node[right,font=\scriptsize,layergrey] at (0.52,0.125) {$n=8$};
  \end{scope}
\end{tikzpicture}
\caption{Apexes that move and apexes that shrink. Left: $g_n$ has constant height
$1$, base $1/n$, and apex at $1-1/n$; the height does not drop,
so $\|g_n\|_\infty=1$ even though $g_n(x)\to0$ at every point. Right: $h_n$ has
height $1/n$ with base $1$ in the same place; its whole graph fits inside the
band of height $1/n$, so it converges uniformly.}
\label{fig:puncak}
\end{figure}

\begin{lemma}
$f_n\to f$ uniformly on $E$ if and only if $\|f_n-f\|_\infty\to0$.
\end{lemma}

\begin{proof}
Uniform convergence means: for every $\varepsilon>0$ there is $N$ such that
for $n\ge N$ we have $|f_n(x)-f(x)|\le\varepsilon$ for all $x\in E$,
that is, $\sup_x|f_n(x)-f(x)|\le\varepsilon$. That is precisely the statement that the
sequence of real numbers $\|f_n-f\|_\infty$ tends to zero. The only technical
difference is that the definition uses $<\varepsilon$ whereas the norm uses
$\le\varepsilon$; the two are equivalent because $\varepsilon$ is arbitrary.
\end{proof}

That lemma turns a question about functions into a question about a single
sequence of numbers: to check uniformity one supremum suffices, to reject it
it suffices to show that this supremum does not tend to zero, as with $x^n$.

Uniformity has an exact picture, namely a band. Take $\varepsilon=0.1$. If
the whole graph of $f_n$ for $n\ge N$ lies between $f-\varepsilon$ and
$f+\varepsilon$, then that single $N$ works for all points. What is checked is
not one point, but whether any part of the graph escapes the band
of width $0.2$. For $g_n$ that band is always pierced by the apex; the
hoped-for situation is visible in Figure~\ref{fig:pita}.

\begin{figure}[htbp]
\centering
\begin{tikzpicture}[x=2.55cm,y=4.0cm,font=\small]
  \draw[->,layergrey] (0,0)--(4.5,0) node[right,font=\scriptsize]{$x$};
  \draw[->,layergrey] (0,0)--(0,1.18) node[above,font=\scriptsize]{$y$};
  \fill[layergrey!12] plot[domain=0:4,samples=90] (\x,{0.55+0.22*sin(2*\x r)+0.1})
    -- plot[domain=4:0,samples=90] (\x,{0.55+0.22*sin(2*\x r)-0.1}) -- cycle;
  \draw[dashed,layergrey] plot[domain=0:4,samples=90] (\x,{0.55+0.22*sin(2*\x r)+0.1});
  \draw[dashed,layergrey] plot[domain=0:4,samples=90] (\x,{0.55+0.22*sin(2*\x r)-0.1});
  \draw[thick,accent] plot[domain=0:4,samples=90] (\x,{0.55+0.22*sin(2*\x r)});
  \draw[thick,layerindigo] plot[domain=0:4,samples=180]
    (\x,{0.55+0.22*sin(2*\x r)+0.06*cos(9*\x r)});
  \draw[<->,layergrey] (3.78,{0.55+0.22*sin(2*3.78 r)-0.1})
    -- (3.78,{0.55+0.22*sin(2*3.78 r)+0.1});
  \node[right,font=\scriptsize,layergrey] at (3.8,{0.55+0.22*sin(2*3.78 r)}) {$0.2$};
  \node[accent,font=\scriptsize,above left] at (0.42,{0.55+0.22*sin(2*0.42 r)}) {$f$};
  \node[layerindigo,font=\scriptsize,below right] at (0.95,{0.55+0.22*sin(2*0.95 r)+0.06*cos(9*0.95 r)}) {$f_n$};
\end{tikzpicture}
\caption{The uniform band of width $0.2$ around $f$, that is between $f-0.1$
and $f+0.1$. The wavy graph of $f_n$ still lies entirely inside the
band, so $\|f_n-f\|_\infty\le0.1$ and one and the same $N$ works for all
points at once.}
\label{fig:pita}
\end{figure}

\begin{tangan}
Take a two-meter rope, tape, and two people. Stretch the rope straight on the
floor, then tape five marks at equal spacing. The first person holds the left end
and lifts it by just $5$ cm. Watch: only that end rises, the other four marks
still touch the floor. That is pointwise convergence.

Now both people lift both ends together on the same count: $5$ cm, then $3$ cm,
then $1$ cm, then half a centimeter. Watch the whole length of the rope. No point is
left on the floor, and the greatest distance from rope to floor is the same at all
points, because both ends rise together. That is uniform convergence.

Finally, the first person lifts his end while walking toward the middle, while
the second lowers hers. The rope stays tilted, and there is always some part
$5$ cm above the floor: every point eventually comes down, but the greatest distance
never shrinks. That is $g_n$ in Figure~\ref{fig:puncak}.
\end{tangan}

\begin{lemma}[Uniform Cauchy criterion]
The sequence $f_n$ converges uniformly on $E$ if and only if for every
$\varepsilon>0$ there is $N$ such that $n,m\ge N$ implies
$|f_n(x)-f_m(x)|<\varepsilon$ for all $x\in E$. This last condition is called the
uniform Cauchy criterion\index{uniform Cauchy criterion}.
\end{lemma}

\begin{proof}
$(\Rightarrow)$ If $f_n\to f$ uniformly, take $N$ from the definition of convergence
for $\varepsilon/2$. For $n,m\ge N$ and every $x$,
\[
  |f_n(x)-f_m(x)|\le|f_n(x)-f(x)|+|f(x)-f_m(x)|<\varepsilon .
\]
$(\Leftarrow)$ For every fixed $x$, the sequence $f_n(x)$ is a Cauchy sequence
in $\mathbb{R}$, so it converges to a number that we call $f(x)$. This
defines a function $f$ point by point. To obtain
uniformity, fix $n\ge N$ in the Cauchy condition and let
$m\to\infty$. Since $f_m(x)\to f(x)$, we get $|f_n(x)-f(x)|\le\varepsilon$
for all $x\in E$ at once, and the $N$ used does not depend on $x$.
\end{proof}

The proof of the second direction uses only the completeness of $\mathbb{R}$: the space of
functions is complete with respect to the supremum norm, and that completeness is what lets a
series of functions be summed without fear that the limit leaves the class of functions we
use.

\begin{catatan}
Uniform convergence is a property of a pair: a sequence and a set. The sentence
``$x^n$ converges uniformly'' is incomplete without naming the set. On
$[0,1]$ the answer is no, on $[0,0.99]$ yes: the supremum norm of the error is
$0.99^n$, namely $0.37$ at $n=100$ and $4.3\times10^{-5}$ at $n=1000$.
Therefore every claim of uniformity in this chapter names its set.
\end{catatan}

\section{The Weierstrass test: uniformity from a single series of numbers}
\tingkat{intermediate}

A series of functions is a series whose terms are functions. Its partial sums
$S_N=\sum_{n=1}^{N}f_n$ are functions, so a series of functions is nothing but a
particular sequence of functions, and all the definitions of the previous section apply
as soon as we replace $f_n$ by $S_N$.

\begin{definisi}
The series $\sum_{n\ge1}f_n$ \emph{converges pointwise} on $E$ if for
every $x\in E$ the series of numbers $\sum_n f_n(x)$ converges. That series
\emph{converges uniformly}\index{series of functions} on $E$ if its sequence of partial
sums converges uniformly on $E$.
\end{definisi}

The following test is the cheapest way to obtain uniformity, and almost every
example in this book uses it. The idea: if each term is bounded by a number
whose series converges, then the tail of the series of functions is bounded by the tail of
that series of numbers, and the numerical tail knows nothing about $x$.

\begin{teorema}[Weierstrass $M$-test]\index{Weierstrass test}
Suppose $f_n$ is defined on $E$ and there are numbers $M_n\ge0$ with
$|f_n(x)|\le M_n$ for all $x\in E$ and all $n$. If $\sum M_n$
converges, then $\sum f_n$ converges uniformly on $E$. If $E$ is a closed
bounded interval and every $f_n$ is continuous, then the sum is continuous on $E$.
\end{teorema}

\begin{proof}
Write $S_N=\sum_{i=1}^{N}f_i$. For $n>m$ and every $x\in E$,
\[
  |S_n(x)-S_m(x)|\le\sum_{i=m+1}^{n}|f_i(x)|\le\sum_{i=m+1}^{n}M_i .
\]
Since $\sum M_n$ converges, the tail $\sum_{i=m+1}^{n}M_i$ can be made smaller
than $\varepsilon$ provided $m,n$ are large enough, and that choice does not
mention $x$. The uniform Cauchy criterion closes the proof of uniformity. For
continuity, each $S_N$ is continuous as a finite sum of continuous functions, and the
uniform limit of continuous functions is continuous; that last statement is proved in
the next section.
\end{proof}

\begin{contoh}
On $E=[-1,1]$ the series $\sum_{n\ge1}x^n/n^2$ converges uniformly, since
$|x^n/n^2|\le1/n^2$ and $\sum1/n^2=\pi^2/6$ converges. Its sum is continuous on
$[-1,1]$, and its value at $x=1$ is $1.6449340668$. The sum of two million terms
gives $1.6449335668$; the difference is $5.0\times10^{-7}$, proportional to the tail
bound $1/N=5\times10^{-7}$. At $x=-1$ the series converges to $-\pi^2/6$, since
the terms alternate in sign with decreasing absolute value tending to zero.
\end{contoh}

\begin{contoh}
On $E=\mathbb{R}$ the series $\sum_{n\ge0}2^{-n}\sin(2^n x)$ converges uniformly,
since $|2^{-n}\sin(2^n x)|\le2^{-n}$ and $\sum2^{-n}=1$. Every term is continuous,
so the sum is continuous at every point. What is surprising is that the sum is not
differentiable at any point: the classical result of Weierstrass (1872), extended
by Hardy up to ratio $1$, which we quote without proof. Uniformity guarantees
continuity, but it guarantees nothing about derivatives.
\end{contoh}

\begin{catatan}
The $M$-test gives a sufficient condition, not a necessary one, and its failure is not a proof
of anything. The series $\sum_{n\ge1}(-1)^n(1-x)x^n$ on $[0,1]$ converges uniformly:
the remainder after the $N$-th term is bounded by $(1-x)x^{N+1}$, whose maximum on $[0,1]$
is $\tfrac{1}{N+2}\left(\tfrac{N+1}{N+2}\right)^{N+1}\to0$, namely
$3.20\times10^{-2}$ at $N=10$ and $3.67\times10^{-4}$ at $N=1000$. But
$\sup_{[0,1]}|f_n|=n^n/(n+1)^{n+1}$, proportional to $1/(en)$, so that
$\sum_n\sup|f_n|$ diverges: $0.88$ up to $N=10$ and $2.53$ up to $N=1000$.
There is no convergent comparison series, so the $M$-test does not apply even though
the uniformity is genuine.
\end{catatan}

\begin{untukanak}
There are two ways to work through a list of chores. The first: finish all the chores
one by one, then add up the total. The second: keep a running total while
working, then finish the rest. For ordinary addition both ways give the
same result, and we are used to swapping them without thinking.

But some tasks do not like being swapped. Imagine you measure the length of a
pole's shadow every hour for a day, compute the average, then
look for the hour when that average is longest. Reverse the order: first find the
longest hour on each day, then average that hour over the month. The two
results can differ widely, not because the arithmetic is wrong, but because
``longest'' and ``average'' genuinely do not commute.

The same thing happens with the formulas used to add up infinitely
many terms, and then integrate or differentiate them. Sometimes the order may be
swapped and the result is exactly the same; sometimes not, and the difference is large. In this chapter
we shall know the warning signs before the result goes wrong.
\end{untukanak}

\section{When limits may be exchanged}
\tingkat{intermediate}

The three operations we habitually use on functions are taking a limit at
a point, integrating, and differentiating. Each of the three is itself a limit:
continuity at $x_0$ means $\lim_{x\to x_0}f(x)=f(x_0)$, the integral is the
limit of Riemann sums, the derivative is the limit of a difference quotient. So moving the
limit of a sequence past one of those operations means exchanging two limits, and
that exchange is not free: its price is one hypothesis, namely
uniformity\index{exchange of limits}.

\begin{teorema}
If $f_n$ is continuous on $E$ and $f_n\to f$ uniformly on $E$, then $f$ is continuous
on $E$.
\end{teorema}

\begin{proof}
Take $x_0\in E$ and $\varepsilon>0$. Since the convergence is uniform, there is $n$
with $\|f_n-f\|_\infty<\varepsilon/3$. Since $f_n$ is continuous at $x_0$, there is
$\delta>0$ such that $|x-x_0|<\delta$ implies $|f_n(x)-f_n(x_0)|<\varepsilon/3$.
Then for $x\in E$ with $|x-x_0|<\delta$,
\[
  |f(x)-f(x_0)|\le|f(x)-f_n(x)|+|f_n(x)-f_n(x_0)|+|f_n(x_0)-f(x_0)|
  <\varepsilon .
\]
\end{proof}

Note the logical order: $n$ is chosen first, then $\delta$, and
$\delta$ may depend on $n$. What is not allowed is choosing a $\delta$
that must work for all $n$ at once.

\begin{tantangan}
The series $\sum_{n\ge0}\bigl(x^{2n}-x^{2n+1}\bigr)$ has continuous terms on $[0,1]$
and converges at every one of its points. Determine its sum, then show that this sum
is not continuous at $x=1$. Why does this result not contradict the theorem above?
A correct answer must name one number that may depend on $\varepsilon$ but must not
depend on $x$.

Hint 1: the partial sums have a closed form, and that form contains
$x^{2N+2}$.

Hint 2: for $x<1$ the sum is $1/(1+x)$, whereas at $x=1$ every term
equals zero; so the value jumps from $1/2$ to $0$.

If you get stuck, go back to the section on the diagonal test and ask what happens at
$x_N=1-1/N$.
\end{tantangan}

\begin{teorema}[Term-by-term integration]
If $f_n$ is Riemann integrable on $[a,b]$ and $f_n\to f$ uniformly on
$[a,b]$, then $f$ is integrable on $[a,b]$ and
\[
  \int_a^b f=\lim_{n\to\infty}\int_a^b f_n .
\]
\end{teorema}

\begin{proof}
For integrability, take $\varepsilon>0$ and choose $n$ with
$\|f_n-f\|_\infty<\varepsilon$. Since $f_n$ is integrable, there is a partition
$\mathcal{P}$ with $U(\mathcal{P},f_n)-L(\mathcal{P},f_n)<\varepsilon(b-a)$.
On each subinterval, $M_i(f)\le M_i(f_n)+\varepsilon$ and
$m_i(f)\ge m_i(f_n)-\varepsilon$, so that
$U(\mathcal{P},f)\le U(\mathcal{P},f_n)+\varepsilon(b-a)$ and
$L(\mathcal{P},f)\ge L(\mathcal{P},f_n)-\varepsilon(b-a)$. Those three inequalities
give $U(\mathcal{P},f)-L(\mathcal{P},f)<3\varepsilon(b-a)$, and the integrability
criterion of Chapter 13 closes the argument. For the value of the integral, it suffices
to note that
\[
  \left|\int_a^b f_n-\int_a^b f\right|
  =\left|\int_a^b (f_n-f)\right|
  \le\int_a^b|f_n-f|\le(b-a)\,\|f_n-f\|_\infty ,
\]
and the right-hand side tends to zero.
\end{proof}

\begin{akibat}[Integration of series]
If $\sum_{n\ge1}f_n$ converges uniformly on $[a,b]$ and every $f_n$ is
integrable there, then
\[
  \int_a^b\sum_{n=1}^{\infty}f_n=\sum_{n=1}^{\infty}\int_a^b f_n .
\]
\end{akibat}

This corollary is applied to the partial sums, with uniformity usually
checked through the $M$-test. For example, the series $\sum_{n\ge1}x^n/n^2$ converges uniformly
on $[0,1]$, so that
\[
  \int_0^1\sum_{n=1}^{\infty}\frac{x^n}{n^2}\,dx
  =\sum_{n=1}^{\infty}\frac{1}{n^2(n+1)}=\frac{\pi^2}{6}-1=0.6449340668 .
\]
That last number is also obtained by summing the first two million terms of the series on
the right-hand side, which gives $0.6449340668$ at machine precision.

\begin{contoh}[Failure without uniformity]
Suppose $f_n(x)=nx(1-x^2)^n$ on $[0,1]$. For every $x\in(0,1]$ the factor
$(1-x^2)^n$ decreases geometrically, so $n(1-x^2)^n\to0$ and $f_n(x)\to0$;
at $x=0$ every $f_n$ equals $0$. So the pointwise limit is the zero function.
Its integral can be computed exactly by the substitution $u=1-x^2$:
\[
  \int_0^1 nx(1-x^2)^n\,dx=\frac{n}{2}\int_0^1 u^n\,du=\frac{n}{2(n+1)} .
\]
That sequence tends to $1/2$, not to $0$. Numerical quadrature gives
$0.4545454545$ for $n=10$ and $0.4950495050$ for $n=100$, matching
$n/(2(n+1))$ to the last digit. So $\lim_n\int_0^1f_n=1/2$ whereas
$\int_0^1\lim_nf_n=0$.
\end{contoh}

That failure has a geometric shape. The function $f_n$ is zero at both ends
and peaks at $x=1/\sqrt{2n+1}$, with height $1.3396685$ at $n=10$ and
$4.2834707$ at $n=100$: the peak grows taller and narrower, so that its
area stays near $1/2$ while its height at every point tends to zero.

\begin{teorema}[Term-by-term differentiation]
Suppose $f_n$ is differentiable on $[a,b]$, the sequence $f_n'$ converges uniformly
on $[a,b]$, and there is one point $x_0\in[a,b]$ where $f_n(x_0)$ converges.
Then $f_n$ converges uniformly to some $f$ that is differentiable on $[a,b]$,
and
\[
  f'(x)=\lim_{n\to\infty}f_n'(x)
\]
for every $x\in[a,b]$.
\end{teorema}

\begin{proof}[Sketch of proof]
The heart of it is the mean value theorem applied to the difference of two terms: for $m,n$ and every
$x,y\in[a,b]$,
\[
  \bigl|(f_n-f_m)(x)-(f_n-f_m)(y)\bigr|\le\|f_n'-f_m'\|_\infty\,|x-y| .
\]
Since $f_n'$ is uniformly Cauchy, the right-hand side can be made small simultaneously for all
pairs $x,y$ once $m,n$ are large. Together with convergence at one point $x_0$,
that inequality makes $f_n$ uniformly Cauchy, so $f_n\to f$ uniformly. For the
derivative, divide by $|x-y|$, let $y\to x$, then exchange the two limits that
appear; the same inequality guarantees that the exchange is legitimate. The technical details are in
Rudin, Theorem 7.17.
\end{proof}

The condition ``there is one point where $f_n(x_0)$ converges'' cannot be dropped: $f_n$
and $f_n+1$ have the same sequence of derivatives, but only one of them
converges at $x_0$. Conversely, uniformity of $f_n$ alone is also not enough to
conclude anything about the derivative.

\begin{contoh}[Derivatives that do not converge]
Suppose $f_n(x)=\sin(nx)/n$ on $[0,\pi]$. Since $\|f_n\|_\infty=1/n$,
the sequence converges uniformly to the zero function; its value is $0.1$ at $n=10$ and
$0.01$ at $n=100$. But $f_n'(x)=\cos(nx)$: at $x=0$ that sequence equals
$1$ for all $n$, whereas at $x=\pi/2$ its values are $0,-1,0,1,0,-1,\ldots$
according to the parity of $n$, which does not converge. So $f_n'$ fails to converge pointwise
anywhere except at special points, let alone uniformly. Differentiation
does not forgive small high-frequency oscillations: $f_n$ shrinks to zero,
but its slope does not.
\end{contoh}

Figure~\ref{fig:turunan} shows three of those derivative curves on $[0,\pi]$. At
$x=\pi/2$ they take the values $1$, $0$, and $-1$: three different values of one
sequence at one point, and that is already enough to say the limit does not exist.

\begin{figure}[htbp]
\centering
\begin{tikzpicture}[x=3.4cm,y=1.5cm,font=\small]
  \draw[->,layergrey] (0,0)--(3.55,0) node[right,font=\scriptsize]{$x$};
  \draw[->,layergrey] (0,-1.28)--(0,1.28) node[above,font=\scriptsize]{$y$};
  \draw[dashed,layergrey] (1.571,-1.25)--(1.571,1.25);
  \node[below,font=\scriptsize,layergrey] at (1.571,-1.25) {$\pi/2$};
  \draw[thick,layerindigo] plot[domain=0:3.1416,samples=160] (\x,{cos(4*\x r)});
  \draw[thick,layerteal] plot[domain=0:3.1416,samples=160] (\x,{cos(5*\x r)});
  \draw[thick,layerorange] plot[domain=0:3.1416,samples=160] (\x,{cos(6*\x r)});
  \fill[ink] (1.571,1) circle (1.5pt);
  \fill[ink] (1.571,0) circle (1.5pt) node[right,font=\scriptsize] {$0$};
  \fill[ink] (1.571,-1) circle (1.5pt) node[right,font=\scriptsize] {$-1$};
  \node[right,font=\scriptsize] at (1.585,1.06) {$1$};
\end{tikzpicture}
\caption{The derivatives $f_n'=\cos(nx)$ on $[0,\pi]$ for $n=4$ (indigo), $n=5$
(teal), and $n=6$ (orange). At $x=\pi/2$ they take the values $1$, $0$, and $-1$: three
different values of one sequence at one point, so the limit of the derivatives does not exist,
even though $f_n=\sin(nx)/n$ converges uniformly to zero with amplitude $1/4$,
$1/5$, and $1/6$.}
\label{fig:turunan}
\end{figure}

\begin{catatan}
In machine learning, the risk function is often written as an expectation over the
data distribution, $\mathbb{E}_{(x,y)\sim D}[\ell(f(x),y)]$, and then differentiated
with respect to the model parameters. If the distribution is finite discrete, as on
the training set, that expectation is a finite sum and exchanging the derivative with
the sum is always legitimate, with no conditions. If it is continuous, the expectation is an integral, and
moving the derivative inside it demands a domination condition more delicate
than uniformity; the tool is the dominated convergence theorem of measure
theory, which is not discussed here. ``Differentiate inside the integral'' is an operation
with a license, not an operation that is always correct.
\end{catatan}

\section{Power series and the radius of convergence}
\tingkat{advanced}

A power series is a series of functions of the most rigid shape: the $n$-th term is
a fixed coefficient times $x^n$. That rigidity bears fruit. The set on which the series
converges is always an interval symmetric about its center, so that all the
information about where the series is useful is contained in one number.

\begin{definisi}
A series of the form $\sum_{n\ge0}a_nx^n$ with $a_n\in\mathbb{R}$ is called a power
series\index{power series} centered at $0$. The number
\[
  R=\sup\Bigl\{r\ge0:\ \sum_{n\ge0}a_nr^n\ \text{converges}\Bigr\}
\]
is called the radius of convergence\index{radius of convergence} of that series.
\end{definisi}

\kamus{power series}{deret pangkat}{a series whose terms are fixed coefficients times
rising powers of one variable, like the same recipe for every value.}

\kamus{radius of convergence}{jari-jari konvergensi}{the furthest distance from the center
at which that recipe still produces a finite number.}

\begin{teorema}
Suppose $\sum a_nx^n$ is a power series with radius $R$. Then the series
converges absolutely for every $|x|<R$ and diverges for every $|x|>R$.
Moreover
\[
  \frac{1}{R}=\limsup_{n\to\infty}|a_n|^{1/n},
\]
with the conventions $1/0=\infty$ and $1/\infty=0$.
\end{teorema}

\begin{proof}
Fix $x\ne0$ and write $|a_nx^n|^{1/n}=|x|\cdot|a_n|^{1/n}$. By the root test
of Chapter 5, the series $\sum a_nx^n$ converges absolutely if
$\limsup_n|a_nx^n|^{1/n}<1$, that is if $|x|<1/\limsup_n|a_n|^{1/n}$, and
diverges if that limit superior is greater than $1$, since its terms do not
tend to zero. Since $x$ is arbitrary, the boundary between the two regimes is exactly
$R=1/\limsup_n|a_n|^{1/n}$, and the set of convergence lies inside $[-R,R]$.
\end{proof}

At $|x|=R$ the root test gives no verdict, and indeed there is no general
verdict: $\sum x^n/n^2$ converges at both ends, $\sum x^n/n$ only at $x=-1$, and
$\sum x^n$ diverges at both.

\begin{akibat}
If $a_n\ne0$ for $n$ large enough and $|a_{n+1}/a_n|\to L$, then $R=1/L$
with the same conventions as above.
\end{akibat}

The reason is the chain of inequalities
\[
  \liminf|a_{n+1}/a_n|\le\liminf|a_n|^{1/n}\le\limsup|a_n|^{1/n}
  \le\limsup|a_{n+1}/a_n|,
\]
which is proved as an exercise in Chapter 5. If the ratio has a limit, all four of those
numbers are equal, so the ratio formula and the root formula give the same $R$. In
hand computation, the ratio formula is almost always faster.

Now the part that makes power series useful: they may be differentiated and
integrated term by term, and the result holds on the whole interval of
convergence, not only on a small piece of it.

\begin{teorema}[Term-by-term differentiation and integration]
Suppose $\sum a_nx^n$ has radius $R>0$. Then for every $x$ with
$|x|<R$,
\[
  \frac{d}{dx}\sum_{n\ge0}a_nx^n=\sum_{n\ge1}n\,a_nx^{n-1},
  \qquad
  \int_0^x\sum_{n\ge0}a_nt^n\,dt=\sum_{n\ge0}\frac{a_n}{n+1}x^{n+1},
\]
and both new series also have radius $R$.
\end{teorema}

\begin{proof}
Take $r$ with $0<r<R$ and choose $\rho$ with $r<\rho<R$. The series
$\sum|a_n|\rho^n$ converges, so there is $C$ with $|a_n|\rho^n\le C$ for all
$n$. For $|x|\le r$,
\[
  |n\,a_nx^{n-1}|\le n\,\frac{C}{\rho}\left(\frac{r}{\rho}\right)^{n-1},
\]
and the series $\sum_n n(r/\rho)^{n-1}$ converges since $r/\rho<1$. The $M$-test
therefore gives uniformity of the differentiated series on $[-r,r]$, and the term-by-term
differentiation theorem gives equality there. Every point $|x|<R$ lies in some
$[-r,r]$ with $r<R$, so the equality holds on all of $(-R,R)$. For the
integral, the same route is used on $\sum|a_n|r^{n-1}$, which also converges.
The radius does not change because the coefficients of the differentiated series are $na_n$, and
\[
  \limsup_n(n|a_n|)^{1/n}
  =\limsup_n|a_n|^{1/n}\cdot\lim_n n^{1/n}
  =\limsup_n|a_n|^{1/n},
\]
since $n^{1/n}\to1$. The same argument applies to the coefficients $a_n/(n+1)$,
because $(n+1)^{-1/n}\to1$.
\end{proof}

That proof uses the whole content of the chapter: the $M$-test for uniformity, the differentiation
theorem to move the limit, and the radius formula to make sure the new series does not
shrink. Uniformity is claimed only on $[-r,r]$, $r<R$; in the opening example,
$\ln(1+x)$ is uniform on every $[-r,r]$ with $r<1$ but not on $(-1,1)$.

\begin{contoh}
For $|x|<1$, the geometric series gives $\sum_{n\ge0}x^n=1/(1-x)$. Differentiating
term by term,
\[
  \sum_{n\ge1}n\,x^{n-1}=\frac{1}{(1-x)^2},
\]
and at $x=1/2$ the right-hand side is $4$. The sum of the first two hundred terms of the series on the
left gives $4.000000$ at machine precision. Integrating term by term from
$0$ to $x$ gives
\[
  \sum_{n\ge1}\frac{(-1)^{n+1}x^n}{n}=\ln(1+x),\qquad |x|<1,
\]
which is the series from the chapter opening. Both results are not numerical coincidences: both are
consequences of the theorem above, because $1/2$ and the $x$ used lie inside
the radius.
\end{contoh}

With that theorem, the opening question of the chapter is answered. The series for $\ln(1+x)$ converges
uniformly on every $[-r,r]$ with $r<1$, since $|(-1)^{n+1}x^n/n|\le r^n/n$
and $\sum r^n/n$ converges. On $(-1,1)$ uniformity fails, and the price shows
in the number of terms: $16$ at $x=0.5$, $5248$ at $x=0.999$. At the endpoint
$x=1$ the series still converges, since it becomes an alternating harmonic series. That is
not a consequence of the radius but of Abel's theorem: if $\sum a_nR^n$ converges,
then $\sum a_nx^n$ converges uniformly on $[0,R]$. We quote it; the proof is in
Rudin, chapter ``Some Special Functions'', Theorem 8.2.

One practical warning: uniform bounds are designed for the worst case. For the
series $\ln(1+x)$ with $N=12$ and $r=0.5$, the bound $r^{13}/13=9.4\times10^{-6}$
is only three times the actual error at $x=0.5$, but at $x=0.1$ it is loose by about
ten thousand times.

\begin{resep}
\begin{enumerate}
\item Write down the difference, $M_N=\sup_{x\in E}|f_N(x)-f(x)|$, or estimate it with
      the value at the worst point.
\item Find the point $x_N$ that maximizes that difference; usually a peak, or
      a point chasing the edge of the interval.
\item Compute $M_N$ as an expression containing only $N$, then check
      $M_N\to0$.
\end{enumerate}
A warning about order: find $x_N$ first for each $N$, then compute its
supremum. Reversing that order, that is choosing $N$ first and then looking for one $x$ that
represents all points, yields a wrong verdict of uniformity for $g_n$ in
Figure~\ref{fig:puncak}, whose apex actually shifts together with $N$.
\end{resep}

\section{Worked numerical examples}
\tingkat{intermediate}

The three groups of numbers below are recomputed in the computational lab. The first
measures the failure of uniformity of $x^n$ on $[0,1]$.

\begin{center}
\small
\begin{tabular}{rrrrr}
\hline
$n$ & $\sup_{[0,1]}x^n$ & $(1-\tfrac1n)^n$ & $(1-\tfrac1{\sqrt n})^n$ & $0.9^{\,n}$\\
\hline
$10$    & $1$ & $0.3486784401$ & $2.234\times10^{-2}$ & $0.3486784401$\\
$100$   & $1$ & $0.3660323413$ & $2.656\times10^{-5}$ & $2.656\times10^{-5}$\\
$1000$  & $1$ & $0.3676954248$ & $1.108\times10^{-14}$ & $1.748\times10^{-46}$\\
$10000$ & $1$ & $0.3678610464$ & $2.249\times10^{-44}$ & $<10^{-450}$\\
\hline
\end{tabular}
\end{center}

The second and fifth columns are two different rates. At the fixed point $x=0.9$
the error is $0.9^n$, falling by a factor of $0.9$ per term: $132$ terms are needed to
reach $10^{-6}$, whereas at $x=0.5$ $20$ terms suffice and at $x=0.99$ it takes
$1375$. Every point has its own exponential rate, with base $x$. The second
column does not move: the uniform error on $[0,1]$ stays $1$ for all $n$, so
its rate of decrease is zero. What is being compared is therefore not two numbers,
but something that tends to zero and something that does not.

The second group of numbers is the error of approximating $e^{x}$ by a Taylor series at
$x=1$. The remainder of the series is bounded by $\frac{e}{(N+1)!}$ for $x\in[0,1]$, since
the $(N+1)$-st derivative of $e^x$ is $e^x\le e$.

\begin{center}
\small
\begin{tabular}{rrrr}
\hline
$N$ & $S_N(1)=\sum_{k=0}^{N}\frac{1}{k!}$ & $e-S_N(1)$ & bound $\frac{e}{(N+1)!}$\\
\hline
$5$  & $2.716666666667$ & $1.615\times10^{-3}$ & $3.775\times10^{-3}$\\
$10$ & $2.718281801146$ & $2.731\times10^{-8}$ & $6.810\times10^{-8}$\\
$15$ & $2.718281828459$ & $5.063\times10^{-14}$ & $1.299\times10^{-13}$\\
$20$ & $2.718281828459$ & $0$ (below precision) & $5.320\times10^{-20}$\\
\hline
\end{tabular}
\end{center}

The bound is always correct and larger than the actual error, by a factor of
about $2.3$ to $2.6$: the remainder is roughly
$\frac{1}{(N+1)!}\left(1+\frac{1}{N+2}+\cdots\right)$, and $e$ appears because
its terms are summed at $x=1$. At $N=20$ the actual error is no longer measurable,
since $5.3\times10^{-20}$ is far below one unit in the last place of a quantity of
size $2.718$, namely about $4.4\times10^{-16}$.

The third group of numbers is three radii of convergence, computed in two
different ways at term $1000$.

\begin{center}
\small
\begin{tabular}{llrr}
\hline
series & $R$ from the formula & $|a_{1000}/a_{1001}|$ & $|a_{1000}|^{1/1000}$\\
\hline
$\sum x^n/n^2$          & $1$              & $1.0020010000$ & $0.9862794856$\\
$\sum \frac{n!}{n^n}x^n$ & $e=2.7182818285$ & $2.7169239322$ & $0.3694916635$\\
$\sum \frac{2^n}{n}x^n$  & $0.5$            & $0.5005000000$ & $1.9862320968$\\
\hline
\end{tabular}
\end{center}

The third column estimates $R$ from the ratio of two consecutive terms, the fourth from
the reciprocal of the root estimate. For the second series, $|a_n/a_{n+1}|=(1+1/n)^n$ rises to
$e$; at $n=1000$ its value is $2.7169239$, off by $0.0014$. For the first series,
$(1+1/n)^2$ decreases to $1$ from above, so the estimate is $1.0020$, and that error
$0.0020$ is a sign that we are only at term $1000$. The fourth column approaches
its limit more slowly, since $|a_n|^{1/n}$ approaches $1$ at rate $|\ln a_n|/n$; at
$n=1000$ the difference is $0.0016$, proportional to $\ln(2\pi n)/(2n)$ times $1/e$.

\kertas{Summary of numbers. The uniform error of $x^n$ on $[0,1]$ is $1$ for every
$n$, while the error at a fixed point $x$ is $x^n$: $132$ terms for $10^{-6}$
at $x=0.9$, $1375$ terms at $x=0.99$. The series for $e^x$ at $x=1$ reaches an error of
$2.7\times10^{-8}$ with $10$ terms, below the bound $6.8\times10^{-8}$.
The radii of the three series are $1$, $e$, and $0.5$; the integral
$\int_0^1 nx(1-x^2)^n dx=n/(2(n+1))$ equals $0.4545$ at $n=10$, while
the limit of the functions is zero.}

\section{Computational lab}
\tingkat{advanced}

The following code runs as it stands on Python 3.14.4 with NumPy 2.5.2 and
SciPy 1.18.0. It recomputes the three groups of numbers above, plus
quadrature for the counterexample on integration.

\begin{lab}
\begin{verbatim}
import math
import numpy as np
from scipy.integrate import quad

print("n      sup[0,1]   (1-1/n)^n   (1-1/sqrt n)^n     0.9^n")
for n in (10, 100, 1000, 10000):
    a = (1 - 1/n)**n
    b = (1 - 1/math.sqrt(n))**n
    print(f"{n:6d} {1.0:9.6f} {a:11.6f} {b:17.3e} {0.9**n:10.3e}")

print()
print("terms of ln(1+x) for error < 1e-6")
for x in (0.5, 0.9, 0.99, 0.999):
    N = 1
    while x**N / N > 1e-6:
        N += 1
    s = sum(((-1)**(k+1)) * x**k / k for k in range(1, N + 1))
    print(f"x={x:6.3f}  N={N:5d}  err={abs(math.log(1+x) - s):.3e}")

print()
print("N      S_N              e - S_N        e/(N+1)!")
for N in (5, 10, 15, 20):
    s = math.fsum(1/math.factorial(k) for k in range(N + 1))
    b = math.e / math.factorial(N + 1)
    print(f"{N:3d} {s:.12f} {math.e - s:14.3e} {b:13.3e}")

print()
print("radius from the ratio of terms 1000 and 1001 (log-space)")
la = lambda n: -2 * math.log(n)
lb = lambda n: math.lgamma(n + 1) - n * math.log(n)
lc = lambda n: n * math.log(2) - math.log(n)
for name, g in (("a", la), ("b", lb), ("c", lc)):
    R = math.exp(g(1000) - g(1001))
    H = math.exp(g(1000) / 1000)
    print(f"series {name}: R = {R!r}   |a_n|^(1/n) = {H!r}")
print("e =", repr(math.e))

print()
print("f_n(x) = n x (1-x^2)^n on [0,1]")
for n in (10, 100):
    val, _ = quad(lambda x, n=n: n * x * (1 - x*x)**n, 0, 1)
    print(f"n={n:4d}  integral={val!r}  n/(2(n+1))={n/(2*(n+1))!r}")
\end{verbatim}
\end{lab}

The output obtained:

\begin{verbatim}
n      sup[0,1]   (1-1/n)^n   (1-1/sqrt n)^n     0.9^n
    10  1.000000    0.348678         2.234e-02  3.487e-01
   100  1.000000    0.366032         2.656e-05  2.656e-05
  1000  1.000000    0.367695         1.108e-14  1.748e-46
 10000  1.000000    0.367861         2.249e-44  0.000e+00

terms of ln(1+x) for error < 1e-6
x= 0.500  N=   16  err=3.049e-07
x= 0.900  N=   89  err=4.478e-07
x= 0.990  N=  721  err=4.915e-07
x= 0.999  N= 5248  err=4.993e-07

N      S_N              e - S_N        e/(N+1)!
  5 2.716666666667      1.615e-03     3.775e-03
 10 2.718281801146      2.731e-08     6.810e-08
 15 2.718281828459      5.063e-14     1.299e-13
 20 2.718281828459      0.000e+00     5.320e-20

radius from the ratio of terms 1000 and 1001 (log-space)
series a: R = 1.0020010000000001   |a_n|^(1/n) = 0.9862794856312105
series b: R = 2.7169239322350776   |a_n|^(1/n) = 0.36949166347196
series c: R = 0.5004999999999973   |a_n|^(1/n) = 1.9862320968418674
e = 2.718281828459045

f_n(x) = n x (1-x^2)^n on [0,1]
n=  10  integral=0.4545454545454544  n/(2(n+1))=0.45454545454545453
n= 100  integral=0.4950495049504951  n/(2(n+1))=0.49504950495049505
\end{verbatim}

Four things must be read from this output. First, the second column stays at $1.000000$
for all four values of $n$, while the fifth column is already $3.5\times10^{-1}$ at
$n=10$: the failure of uniformity shows up as a number. Second, the last entry of the fifth
column prints $0.000e+00$ not because the error is zero, but because
$0.9^{10000}=10^{-457.6}$ lies below the smallest normal \texttt{float64} bound: underflow.

Third, the fourth column gives numerical evidence for two limits already
proved: $(1-1/n)^n\to1/e=0.3678794$ and $(1-1/\sqrt n)^n\to0$, which collapses
far faster.

Fourth, SciPy's quadrature gives $0.4545454545454544$ for $n=10$, whereas the
formula $n/(2(n+1))$ gives $0.45454545454545453$; a difference of one unit in the last place
of rounding. Without a closed form, we would have no way to make sure the quadrature is not
measuring something else.

One thing that does not show in the output: every bound in this chapter is an upper bound.
The row $N$ that the code reports is chosen by checking $x^N/N$, so that $N$ is
guaranteed to suffice, not to be minimal; the actual errors are $3.0\times10^{-7}$ to
$5.0\times10^{-7}$, below $10^{-6}$.

\begin{untukdosen}
The teaching order I suggest: start from the failure of the $\ln(1+x)$ library in
the opening, then the pointwise and uniform definitions with $x^n$ as the single
example. Do not rush to the $M$-test. Spend one full session making sure
the students can say precisely why $x^n$ fails to be uniform
and why $(1-1/\sqrt n)^n$ succeeds. After that the $M$-test, then the three exchange
theorems, each with its own counterexample.

Two misconceptions appear most often. First, students think ``every
$f_n$ continuous and $f_n\to f$'' already suffices to conclude that $f$ is continuous; they
forget that what is needed is uniformity, and $x^n$ is the cure. Second,
students think uniformity is inherited downward: if $f_n\to f$ uniformly then
$f_n'\to f'$. The example $\sin(nx)/n$ breaks that in one line; it is best
written on the board before the differentiation theorem appears.

A board activity that works: write four unnamed sequences of functions,
namely $x^n$, $x^n/n$, $nx(1-x^2)^n$, and $\sin(nx)/n$, then ask the class
to decide for each whether it converges pointwise, whether uniformly, and
whether its integral and derivative may be exchanged. The answers form a $4\times4$ table
that can be filled in twenty minutes. Students are ready to go on when they can
explain why a power series is uniform only on $[-r,r]$.
\end{untukdosen}

\section{Summary}
\tingkat{basic}

\begin{ringkasan}
\begin{itemize}
\item Pointwise convergence gives every point its own threshold;
      uniform convergence demands one threshold for all points at once,
      and that difference is equivalent to $\|f_n-f\|_\infty\to0$.
\item The Weierstrass $M$-test concludes uniformity from a single comparison
      series of numbers. It is a sufficient condition, not a necessary one; its failure proves
      nothing.
\item The uniform limit of continuous functions is continuous, and its integral may be exchanged with
      the limit. Differentiation demands uniformity of the sequence of derivatives plus convergence
      at one point.
\item Those three hypotheses cannot be removed: $x^n$, $nx(1-x^2)^n$, and
      $\sin(nx)/n$ each break one of them.
\item A power series converges absolutely inside its radius and diverges outside;
      inside the radius it may be differentiated and integrated term by term
      without changing its radius.
\item Uniformity is always local: $\ln(1+x)$ is uniform on every $[-r,r]$, $r<1$,
      but not on $(-1,1)$, and the price shows in the number of terms, from
      $16$ at $x=0.5$ to $5248$ at $x=0.999$.
\end{itemize}
\end{ringkasan}

\section{Exercises}
\tingkat{intermediate}

\begin{latihan}
\begin{enumerate}
\item Compute in Python, for $n=10,100,1000$, the value
      $\sup_{x\in[0,1]}|x^n|$ and the value of $x^n$ at $x=0.9$. Explain why
      the first column is flat and the second is not.

\item Prove: if $f_n\to f$ uniformly on $E$ and every $f_n$ is bounded on
      $E$, then $f$ is bounded on $E$ and the sequence $\|f_n\|_\infty$ is bounded.
      Show by example that pointwise convergence is not enough.

\item Suppose $f_n(x)=x/(1+nx^2)$ on $\mathbb{R}$. Show that
      $f_n\to0$ pointwise, determine $\|f_n\|_\infty$, and prove that
      the convergence is uniform. What is the smallest $n$ such that
      $\|f_n\|_\infty<0.01$?

\item Judge this claim: true, false, or not enough evidence. ``Since $x^n\to0$
      for every $x\in[0,1)$, the sequence $f_n(x)=x^n$ converges uniformly on
      $[0,1)$.'' State exactly where that reasoning breaks, and give one
      sequence of points that proves it.

\item Compute the error $|e-S_N(1)|$ for $N=5,10,15,20$ and compare it with
      the bound $e/(N+1)!$. What is the smallest $N$ for which that bound is below
      $10^{-16}$? Can the actual error still be measured at that $N$?

\item Prove that the series $\sum_{n\ge0}x^n(1-x)$ converges pointwise
      on $[0,1]$ but not uniformly, by computing its sum first
      and then checking its continuity.

\item Compute the radii of convergence of $\sum x^n/n^2$, $\sum(n!/n^n)x^n$, and
      $\sum 2^nx^n/n$ from the ratio of terms $1000$ and $1001$, then compare
      with the exact values. Which series has the slowest estimate, and
      why?

\item Prove that if $f_n\to f$ uniformly on $E$ and $g$ is bounded on $E$,
      then $f_ng\to fg$ uniformly on $E$. Show by example that the
      boundedness of $g$ cannot be removed.

\item For $f_n(x)=nx(1-x^2)^n$ on $[0,1]$, compute the peak point, the peak
      height, and $\int_0^1f_n$ for $n=10,100$. Compare the peak height
      with $\sqrt{n/(2e)}$ and explain why the area does not tend to zero.

\item Differentiate $\sum_{n\ge0}x^n=1/(1-x)$ term by term and use the result
      to compute $\sum_{n\ge1}n/2^n$. How many terms are needed for the
      partial sum to lie within $10^{-10}$ of the exact value?
\end{enumerate}

\noindent\emph{Hints for the three hardest items.} Item 4: the uniform property
demands one $N$ for all $x$, so it suffices to show that for every $n$
there is an $x$ with $x^n\ge\tfrac12$; take $x=2^{-1/n}$. Item 6: sum the geometric
series first; the result is a function that differs from its limit at one point,
and the uniform limit of continuous functions must be continuous. Item 8: write
$|f_ng-fg|=|f_n-f|\,|g|$ and use $\|g\|_\infty$; for the counterexample
take $g(x)=1/x$ on $(0,1)$ and an $f_n$ that decreases uniformly but not
fast enough.
\end{latihan}

\begin{untukpeneliti}
What is proved with complete proof in this chapter: the equivalence of uniformity with the
supremum norm, the uniform Cauchy criterion, the Weierstrass $M$-test, the three exchange
theorems together with their counterexamples, the radius formula through the limit superior, and
the invariance of the radius under term-by-term differentiation and integration. Two results are
quoted without proof: the non-differentiability of the Weierstrass function (1872, extended
by Hardy 1916) and Abel's theorem. The primary sources are Rudin, \emph{Principles of
Mathematical Analysis}, chapters ``Sequences and Series of Functions'' (7.1--7.17) and
``Some Special Functions'' (8.1--8.2); Abbott, \emph{Understanding Analysis}, chapter
``Sequences and Series of Functions''; Stewart, \emph{Calculus}, the section ``Power
Series''.

What is often popularized beyond its due is the idea that uniformity removes all
problems. It removes three particular problems, on the interval where it holds, and
no more. For Fourier series uniformity is in fact rare, and the exchange of
integral remains legitimate for other reasons. In neural-network approximation, what is used is
convergence on a finite set.

Open question: if $f_n$ is known only through its values on a finite grid, how large
is the error of estimating $\|f_n\|_\infty$ by the maximum on the grid? Without
regularity assumptions there is no general answer, and for narrow peaks a grid that looks
dense can miss them entirely. Limits of validity: all functions here are
real-valued on subsets of $\mathbb{R}$, with 64-bit arithmetic in
Python 3.14.4, NumPy 2.5.2, and SciPy 1.18.0.
\end{untukpeneliti}

\peta{Chapter 14 in Rudin, \emph{Principles of Mathematical Analysis}, chapters
``Sequences and Series of Functions'' (7.1--7.17) and ``Some Special Functions''
(8.1--8.2); in Abbott, \emph{Understanding Analysis}, chapter ``Sequences and
Series of Functions''; in Stewart, \emph{Calculus}, the sections ``Power Series'' and
``Taylor and Maclaurin Series''. Its basis is the second half of the chapter \emph{Integral, Series
of Functions, and Exchange of Limits} of the first edition of this book. This chapter uses the definitions of
limit and the Cauchy criterion from Chapter 4, the root test of Chapter 5, uniform continuity of Chapter 11,
and the integrability criterion of Chapter 13, all of them referred to, not repeated.}
